\documentclass[11pt]{article}

\usepackage{amsmath}
\usepackage{amssymb}
\usepackage{array}
\usepackage{bm}
\usepackage{esvect}
\usepackage{geometry}
\usepackage{makecell}
\usepackage{mathtools}
\usepackage{siunitx}

\usepackage[T1]{fontenc}
\usepackage[utf8]{inputenc}

\geometry{margin=1in}

\DeclareRobustCommand{\myvec}[1]{\vv{\mathbf{#1}}}
\newcommand{\myunitvec}[1]{\bm{\hat{\mathbf{#1}}}}
\newcommand{\myvecsp}[1]{\bm{\mathcal{#1}}}
\newcommand{\mysubvecsp}[1]{\bm{{#1}}}
\newcommand{\myset}[1]{\mathcal{#1}}
\newcommand{\mymat}[1]{#1}
\newcommand{\dotprod}[2]{#1 \cdot #2}
\newcommand{\innerprod}[2]{\left\langle #1, #2\right\rangle}

\title{Diagonalization of Matrices}
\author{\texttt{O5-37-2718281828459}}

\raggedright
\begin{document}

    \maketitle

    \newpage
    \vspace*{\fill}
    \begin{center}
        \LARGE
        \fontfamily{qzc}\selectfont\itshape
        The people who seek the truth are more important than the truth itself \\
    \end{center}
    {
        \Large
        \rightline{--- \emph{6, Reverse 1999}}
    }
    \vspace*{\fill}

    \newpage
    \tableofcontents
    
    \newpage
    \setcounter{section}{-1}
    \section{Preliminaries}
    \subsection{The Scalar Field}
    Unless otherwise specified:
    \begin{itemize}
        \item All matrices in this notes are real matrices
        \item All vectors are real
        \item All vector spaces are $\subseteq \mathbb{R}^n$
    \end{itemize}

    \subsection{Coordinate Spaces}
    We distinguish between an abstract subspace $V$ and the entire $n$-dimensional coordinate universe:
    \begin{itemize}
        \item $\mathbb{R}^n$ represents the \textbf{entire vector space} (the domain and codomain of our transformations).
        \item $\mysubvecsp{V}, \mysubvecsp{E}, \dots$ is used to denote an arbitrary subspace or a solution set (such as an eigenspace) within $\mathbb{R}^n$.
    \end{itemize}

    \newpage
    \section{Change of Basis}
    \subsection*{Quick Review: From Axioms to Coordinates}
    Let $\mysubvecsp{V}$ be a vector space and $\myset{B}$ be any basis for $\mysubvecsp{V}$
    \begin{itemize}
        \item \textbf{Basis Uniqueness:} For any $\myvec{v} \in \mysubvecsp{V}$, the coordinate vector $[\myvec{v}]_\myset{B}$ is unique.
        \item \textbf{Origin Invariance:} All coordinate systems in $\mysubvecsp{V}$ share the same $\myvec{0}$ vector.
        \item \textbf{Linearity:} The transition between $[\myvec{v}]_{\mathcal{B}_1}$ and $[\myvec{v}]_{\mathcal{B}_2}$ is a \textit{linear transformation}.
    \end{itemize}
    \subsection{Core Requirement: Shared Vector Space}
    For the change of basis $[x]_\myset{B} \to [x]_\myset{C}$ to be valid, both sets must be bases for the \textbf{same underlying vector space} $\mysubvecsp{V}$.
    \begin{equation}
        \myset{B}, \myset{C} \subset \mysubvecsp{V} \quad \text{and} \quad \operatorname{Span}(\myset{B}) = \operatorname{Span}(\myset{C}) = \mysubvecsp{V}
    \end{equation}
    This ensures that every vector $\myvec{v} \in \mysubvecsp{V}$ has a unique representation in both coordinate systems.
    \subsection{The Transition Matrix}
    Since change of basis is a linear transformation, it can be expressed using a matrix
    \begin{quote}
        Note: in the later parts, we will use the following notations: \vspace{\abovedisplayshortskip}

        \begin{tabular}{l l @{}}
            $\myvec{v}$ & a vector v in space, in standard basis, or independent on any basis \\[1ex]
            $v_k$ & the k-th component of $\myvec{v}$ \\[1ex]
            $[\myvec{v}]_\myset{B}$ & the vector $\myvec{v}$ expressed in a basis $\myset{B}$ \\[1ex]
            $(v_\myset{B})_k$ & the k-th component of $[\myvec{v}]_\myset{B}$ \\[1ex]
            $\mymat{P}_{\myset{B}_2\leftarrow\myset{B}_1}$ & a transition matrix $\mymat{P}$ that translates a vector from basis $\myset{B}_1$ to basis $\myset{B}_2$
        \end{tabular}
    \end{quote}
    \subsubsection{Transit From Standard Basis}
    Let $\myset{S}$ be the \textbf{standard basis} of $\mathbb{R}^n$, $\myset{B} = \{\myvec{b}_1, \dots, \myvec{b}_n\}$ be a basis for $\mathbb{R}^n$, and $\myvec{v} \in \mathbb{R}^n$\\
    if we want to find $[\myvec{v}]_\myset{S}$ from $[\myvec{v}]_\myset{B}$, we have
    \begin{align*}
    \myvec{v} = [\myvec{v}]_{\myset{S}} = {(v_\myset{B})}_1\myvec{b}_1 + {(v_\myset{B})}_2\myvec{b}_2 + \dots + {(v_\myset{B})}_n\myvec{b}_n 
        & = \bigg[ \myvec{b}_1 \, \myvec{b}_2 \, \dots \, \myvec{b}_n \bigg] 
        \begin{bmatrix} {(v_\myset{B})}_1 \\ {(v_\myset{B})}_2 \\ \vdots \\ {(v_\myset{B})}_n \end{bmatrix} \\
        & = \bigg[ \myvec{b}_1 \, \myvec{b}_2 \, \dots \, \myvec{b}_n \bigg]\,[\myvec{v}]_\myset{B}
    \end{align*}
    Therefore, the transition matrix that translate vectors in basis $\myset{B}$ to standard basis $\myset{S}$, is given by,
    \begin{equation}
        \mymat{P}_{\myset{S}\leftarrow\myset{B}} = \bigg[ \myvec{b}_1 \, \myvec{b}_2 \, \dots \, \myvec{b}_n \bigg]
    \end{equation} 
    and to find expression in $\myset{S}$ of a vector expressed in basis $\myset{B}$, we have:
    \begin{equation}
        \myvec{v} = [\myvec{v}]_{\myset{S}} = \mymat{P}_{\myset{S}\leftarrow\myset{B}} \, [\myvec{v}]_{\myset{B}}
    \end{equation}
        
    \subsubsection{General Formula for Transiting in Two Basis}

    Let $\myset{B} = \{\myvec{b}_1, \dots, \myvec{b}_n\}$ and $\myset{C}$ be two bases for $\mysubvecsp{V}$.
    To find the coordinate vector $[\myvec{v}]_{\myset{C}}$ from $[\myvec{v}]_{\myset{B}}$, we use the transition matrix $\mymat{P}_{\myset{C} \leftarrow \myset{B}}$:

    \begin{equation}
        [\myvec{v}]_{\myset{C}} = \mymat{P}_{\myset{C} \leftarrow \myset{B}}\,[\myvec{v}]_{\myset{B}}
    \end{equation}

    The columns of this matrix are the coordinate vectors of the basis $\myset{B}$ relative to the basis $\myset{C}$:

    \begin{equation}
        \mymat{P}_{\myset{C} \leftarrow \myset{B}} = 
        \begin{bmatrix} 
            [\myvec{b}_1]_{\myset{C}} & [\myvec{b}_2]_{\myset{C}} & \dots & [\myvec{b}_n]_{\myset{C}} 
        \end{bmatrix}
    \end{equation}

    This matrix essentially "encodes" how to map the unit steps of the first grid onto the second grid.

    \newpage
    \subsection{Change of Basis for a Linear Transformation}
    
    Think of an example, in 2d, in the standard basis, if we want to rotate our plane anti-clockwise by $\pi/2$ (or $\ang{90}$), we use a rotation matrix $\mymat{R}$:
    \[ \mymat{R} = \begin{bmatrix}
        0 & -1 \\
        1 & 0
    \end{bmatrix} \]
    what if we want a matrix of transformation of "anti-clockwise $\pi/2$ rotation" that applicable on vectors in another basis $\myset{B}$?\\
    or in other words, what is $[\mymat{R}]_\myset{B}$?
    \vspace{\belowdisplayskip}

    \subsubsection{Transitioning from the Standard Basis}

    Let $[\myvec{u}]_{\myset{B}}$ be a vector expressed in basis $\myset{B}$\\
    we want to find $[\myvec{v}]_{\myset{B}}$, \\
    \quad the resultant vector after any linear transformation $\mymat{L}$, expressed in basis $\myset{B}$ \vspace{\belowdisplayshortskip}

    Since we don't know how to do the transformation directly in basis $\myset{B}$, we \\
    \begin{quote}
        first, transit it into standard basis
        \[ \myvec{u} = \mymat{P}_{\myset{S}\leftarrow\myset{B}}\,[\myvec{u}]_{\myset{B}} \]
        then, we do the transformation to the vector, in our familiar standard basis,
        \[ \myvec{v} = \mymat{L}\myvec{u} \]
        afterall, we transit it back to expressing in basis $\myset{B}$ \\
        \[ [\myvec{v}]_{\myset{B}} = \mymat{P}_{\myset{B}\leftarrow{S}} \myvec{v} \qquad \text{where } \mymat{P}_{\myset{B}\leftarrow{S}} = (\mymat{P}_{\myset{S}\leftarrow{B}})^{-1}\]
    \end{quote}
    Therefore, given a transition matrix $\mymat{P} = \mymat{P}_{\myset{S}\leftarrow\myset{B}}$ \\
    for any linear transformation $\mymat{L} = [\mymat{L}]_{\myset{S}}$,\\ 
    the equation to find $[\myvec{v}]_{\myset{B}}$ from $[\myvec{u}]_{\myset{B}}$:
    \begin{equation}
        [\myvec{v}]_{\myset{B}} = \mymat{P}^{-1} \mymat{L} \mymat{P} [\myvec{u}]_{\myset{B}}
    \end{equation}

    Therefore, we have the formula for $[\mymat{L}]_{\myset{B}}$:
    \begin{equation}
        [\mymat{L}]_{\myset{B}} = \mymat{P}^{-1} \mymat{L} \mymat{P}
    \end{equation}

    \subsubsection{General Formula}
    For any linear transformation $[\mymat{L}]_{\myset{B}}$ in basis $\myset{B}$, \\
    the same transformation that applicable on vectors in basis $\myset{C}$, $[\mymat{L}]_{\myset{C}}$, is given by
    \begin{equation}
        [\mymat{L}]_{\myset{C}} = \mymat{P}_{\myset{C}\leftarrow\myset{B}} \cdot [\mymat{L}]_{\myset{B}} \cdot \mymat{P}_{\myset{B}\leftarrow\myset{C}}
    \end{equation}
    
    \newpage
    \section{Eigenvectors}

    \subsection{Definitions}
    \subsubsection{Eigenvectors and Eigenvalues}
    \textbf{Eigenvector} (or characteristic vector) is a non-zero vector that \\
    has its direction un\large \textbf{changed} (or reversed, or shrunk into a point) in a linear transformation
    \vspace{\belowdisplayshortskip}

    More precisely, an eigenvector $\myvec{v}$ of a linear transformation $\mymat{A}$ is scaled by a constant factor $\lambda$ when the linear transformation is applied to it
    \vspace{\belowdisplayshortskip}

    \noindent
    \begin{minipage}{\linewidth}
        An eigenvector $\myvec{v}$ suits the following equation
        \begin{equation}
            \mymat{A}\myvec{v} = \lambda\myvec{v}, \qquad \myvec{v} \ne \myvec{0}
        \end{equation} 
        \begin{tabbing}
            where \= $\mymat{A}$ is any linear transformation, and\\
            \> the scalar $\lambda$ is the \textbf{eigenvalue} (or characteristic value) of $\mymat{A}$ 
        \end{tabbing}
    \end{minipage}
    \hbox{}

    \noindent \rule{\textwidth}{0.1pt} 
    if we have
    \[ \mymat{A}\,\myvec{x}_1 = \lambda_1 \myvec{x}_1, \qquad \myvec{x}_1 \ne \myvec{0} \]
    \begin{center}
        where $\lambda_1$ is one of the eigenvalues
    \end{center}
    then, $\myvec{x}_1$ is called the eigenvector(s) of $\mymat{A}$ corresponds to $\lambda_1$\\
    the solution set for $\myvec{x}_1$, plus the zero vector, is called the eigenspace of $\lambda_1$, denoted by $\mysubvecsp{E}_{\lambda_1}$
    \noindent \rule{\textwidth}{0.1pt} 

    \subsubsection{Eigenspace}
    For a given eigenvalue $\lambda_k$, the \textbf{eigenspace} $\mysubvecsp{E}_{\lambda_k}$ is a vector space, it is the set of all eigenvectors corresponding to $\lambda_k$, plus the zero vector.
    \begin{equation}
        \mysubvecsp{E}_{\lambda_k} = \{ \myvec{v} \mid \mymat{A}\myvec{v} = \lambda_k\myvec{v} \} = \operatorname{null}(\mymat{A}-\lambda_k \mymat{I})
    \end{equation}

    \newpage
    \subsection{Characteristic Equation}
    The characteristic equation is an equation that is very useful for \\
    finding eigenvalues (and other eigen-related infos) for a square matrix $\mymat{A}$
    \begin{equation}
        \det(\mymat{A}-\lambda \mymat{I}) = 0
    \end{equation}

    \subsubsection{Derivation}
    \begin{align*}
        \mymat{A}\myvec{v} &= \lambda \myvec{v}\\
        \mymat{A}\myvec{v} - \lambda \myvec{v} &= \myvec{0}\\
        (\mymat{A} - \lambda \mymat{I})\myvec{v} &= \myvec{0}
    \end{align*}
    This is a homogeneous linear system\\
    \begin{quote}
        for a homogeneous linear system $\mymat{A}\myvec{x} = \myvec{0}$, \\
        \quad if $\det\mymat{A}\ne0$ \quad there are only the trivial solution ($\myvec{x} = \myvec{0}$) \\
        \quad if $\det\mymat{A}=0$ \quad there are infinite many solutions for $\myvec{x}$
    \end{quote}

    since eigenvectors are non-zero vectors, \\
    we want to find the solution $\myvec{v} \ne \myvec{0}$ for $(\mymat{A} - \lambda \mymat{I})\myvec{v} = \myvec{0}$,\\
    $\therefore$ we look for \fbox{$\det(\mymat{A} - \lambda \mymat{I}) = 0$} \quad (infinite many solutions)\\
    and this is our \textbf{characteristic equation}


    \newpage
    \subsection{Notable Properties}

    \subsubsection{Number of Eigenvalues and Dimensions of Eigenspaces}
    Let $\mymat{A}$ be any real $n \times n$ matrix.
    \begin{itemize}
        \item \textbf{Number of Eigenvalues:} \vspace{1ex} \\
            There are always $n$ eigenvalues when counting \textbf{algebraic multiplicities}. Let them be $\lambda_1, \lambda_2, \dots, \lambda_n$.
            \begin{enumerate}
                \item They can be duplicated (i.e.: $\lambda_i = \lambda_j$).
                \item They can be complex (i.e.: $\lambda = a + bi$).
            \end{enumerate} \vspace{1ex}
            If we restrict our count to \textbf{distinct} and \textbf{real} eigenvalues:
            \begin{enumerate}
                \item $\mymat{A}$ may have \textbf{no} real eigenvalues (the spectrum is purely complex).
                \item The number of real eigenvalues $m$ ranges from $0 \le m \le n$.
            \end{enumerate}

        \item \textbf{Dimension of Eigenspaces:} \vspace{1ex} \\
            Let $m$ be the number of \textbf{distinct real} eigenvalues $\lambda_1, \dots, \lambda_m$ with corresponding eigenspaces $\mysubvecsp{E}_{\lambda_1}, \dots, \mysubvecsp{E}_{\lambda_m}$.
            \begin{itemize}
                \item \textbf{Real:} $\mysubvecsp{E}_{\lambda_k} \subset \mathbb{R}^n$ for all real $\lambda_k$.
                \item \textbf{Disjointness:} Eigenspaces only intersect at the zero vector:
                \begin{equation}
                    \mysubvecsp{E}_{\lambda_i} \cap \mysubvecsp{E}_{\lambda_j} = \{\myvec{0}\} \quad \forall\,i \ne j
                \end{equation}
                \item \textbf{Multiplicity Constraint:} For each $\lambda_k$, the \textbf{geometric multiplicity} ($\dim \mysubvecsp{E}_{\lambda_k}$) satisfies:
                \begin{equation}
                    1 \le \dim \mysubvecsp{E}_{\lambda_k} \le (\text{Algebraic Multiplicity of } \lambda_k)
                \end{equation}
                \item \textbf{Total Dimension:} The sum of dimensions $\sum_{k=1}^m \dim \mysubvecsp{E}_{\lambda_k}$ ranges from $m$ to $n$.
                \item \textbf{Full Coverage:} We have: total dimension $ = n$ if $m = n$. \\ 
                    \begin{quote}
                        This implies that the direct sum of all $\mysubvecsp{E}_{\lambda}$ covers $\mathbb{R}^n$ when we have all $\lambda_k$ being distinct. \\
                        This explains why having all distinct real eigenvalues is a sufficient (but not necessary) condition for a square matrix to be \textbf{diagonalizable}, which we will cover in later chapters.
                    \end{quote}
            \end{itemize}
    \end{itemize}
    
    \newpage
    \subsubsection{Eigenvalue of Inverse Matrix}
    If $\mymat{A}^{-1}$ exists, the eigenvalues of $\mymat{A}^{-1}$ are the reciprocal of eigenvalues of $\mymat{A}$\\
    
    derivation: Assume $\mymat{A}$ is invertable, let $\lambda$ and $\lambda'$ be the eigenvalues of $\mymat{A}$ and $\mymat{A}^{-1}$ respectively\\
    Note that $\lambda \ne 0$ is always true when $\mymat{A}$ is invertible (think of the characteristic equation)
    \begin{align*}
        \mymat{A}\myvec{v} &= \lambda\myvec{v} \\
        \mymat{A}^{-1} \mymat{A} \myvec{v} &= \mymat{A}^{-1} (\lambda \myvec{v}) \\
        \myvec{v} &= \lambda \mymat{A}^{-1} \myvec{v} & \mymat{L}(k\myvec{v}) = k\mymat{L}(\myvec{v}) \text{ for any lin. transf. function }\mymat{L}\\
        (1/\lambda)\myvec{v} &= \mymat{A}^{-1} \myvec{v}\\
        \lambda' &= 1/\lambda & \text{by definition of eigenvalue}
    \end{align*}
    \subsubsection{Product and Sum of Eigenvalues}
    \paragraph{The Determinant Property}~\\
    for any $n \times n$ square matrix $\mymat{A}$ with eigenvalues $\lambda_1, \lambda_2, \dots, \lambda_n$
    \begin{equation}
        \prod_{k=1}^{n} \lambda_k = \det \mymat{A}
    \end{equation}

    \paragraph{Proof of, $\prod_{k=1}^{n} \lambda_k = \det \mymat{A}$}~\vspace{1ex}\\
    Let $\mymat{A}$ be any $n \times n$ square matrix with eigenvalues $\lambda_1, \lambda_2, \dots, \lambda_n$ \\
    Let polynomial $P(\lambda) = \det(\mymat{A} - \lambda\mymat{I})$ \vspace{1ex}\\
    
    Since, the term $\lambda^n$ in $\det(\mymat{A} - \lambda\mymat{I})$ can only be produced from the main diagonal of $\mymat{A}$, \\
    the term of $\lambda^n$ in $P(\lambda)$ is $(-1)^n\lambda^n$ \vspace{1ex} \\

    By Fundamental Theorem of Algebra, \\
    $ P(\lambda) = (-1)^n(\lambda - \lambda_1)(\lambda - \lambda_2)\dots(\lambda - \lambda_n)$ \vspace{1ex} \\

    Since $P(\lambda) = \det(\mymat{A} - \lambda\mymat{I})$ for any $\lambda$, we can substitute $\lambda = 0$, \vspace{1ex} \\
    $P(0) = \det\mymat{A}$ \vspace*{-\abovedisplayskip}
    \begin{flalign*}
        (-1)^n( - \lambda_1)( - \lambda_2)\dots( - \lambda_n) &= \det\mymat{A} &\\
        (-1)^n (-1)^n(\lambda_1 \lambda_2 \dots \lambda_n) &= \det\mymat{A} &\\
        (\lambda_1 \lambda_2 \dots \lambda_n) &= \det\mymat{A} &\\
        \Aboxed{ \prod_{k = 1}^{n} \lambda_k &= \det\mymat{A} }
    \end{flalign*}
    \newpage
    
    \paragraph{The Trace Property}~\\
    for any $n \times n$ square matrix $\mymat{A}$ with eigenvalues $\lambda_1, \lambda_2, \dots, \lambda_n$
    \begin{equation}
        \sum_{k=1}^{n} \lambda_k = \operatorname{tr}(\mymat{A})
    \end{equation}
    where $\operatorname{tr}(\mymat{A})$ is the trace of $\mymat{A}$, which equals to $\sum_{i=1}^{n}a_{ii}$ , the sum of main elements at the diagonal of $\mymat{A}$ 
    
    \paragraph{Proof of, $\sum_{k=1}^{n} \lambda_k = \operatorname{tr}(\mymat{A})$}~\vspace{1ex}\\
    Let $\mymat{A}$ be any $n \times n$ square matrix with eigenvalues $\lambda_1, \lambda_2, \dots, \lambda_n$ \\
    Let polynomial $P(\lambda) = \det(\mymat{A} - \lambda\mymat{I})$ 
    \begin{enumerate}
        \item Same as above, $P(\lambda) = (-1)^n(\lambda - \lambda_1)(\lambda - \lambda_2)\dots(\lambda - \lambda_n)$, \\ 
            notice that:
            \begin{quote}
                coefficient of $\lambda^n = (-1)^n$ \\
                coefficient of $\lambda^{n-1} = (-1)^n (-\lambda_1-\lambda_2-\dots-\lambda_n) = (-1)^{n+1} \sum_{k=1}^{n}\lambda_k$
            \end{quote}
        \item Observe the calculation of $\det(\mymat{A}-\lambda\mymat{I})$, \\
            notice that, the terms of $\lambda^n$ and $\lambda^{n-1}$ comes out only from the main diagonal:
            \[ (a_{11} - \lambda)(a_{22} - \lambda)\dots(a_{nn} - \lambda) \]
            so, 
            \begin{quote}
                coefficient of $\lambda^n = (-1)^n$ \\
                coefficient of $\lambda^{n-1} = (-1)^n (-a_{11}-a_{22}-\dots-a_{nn}) = (-1)^{n+1} \sum_{k=1}^{n}a_{kk}$
            \end{quote}
    \end{enumerate}
    Therefore, we have 
    \begin{gather*}
        (-1)^{n+1} \sum_{k=1}^{n}\lambda_k = \text{ coefficient of }\lambda^{n-1} = (-1)^{n+1} \sum_{k=1}^{n}a_{kk} \\
        \sum_{k=1}^{n}\lambda_k = \sum_{k=1}^{n}a_{kk} \\
        \boxed{\sum_{k=1}^{n}\lambda_k = \operatorname{tr}(\mymat{A})}
    \end{gather*}
    
    
    
    

    \newpage
    \section{Diagonalization}
    \subsection{Diagonalizable Matrices}
    For a $n \times n$ square matrix $\mymat{A}$ that has $m$ distinct eigenvalues $\lambda_1, \lambda_2, \dots, \lambda_m$ to be \textbf{diagonalizable}, the total dimension of its eigenspaces spans the entire vector space
    \begin{equation}
        \mymat{E}_{\lambda_1} \oplus \mymat{E}_{\lambda_2} \oplus \dots \oplus \mymat{E}_{\lambda_m} = \mathbb{R}^n
    \end{equation}
    In practical calculation, this means the sum of the \textbf{geometric multiplicities} must equal the order of the matrix:
    \begin{equation}
        \sum_{k=1}^{m} \dim(\mysubvecsp{E}_{\lambda_k}) = n
    \end{equation}
    Property: \\
    if a square matrix has all the $n$ eigenvalues being distinct $\implies$ it is diagonalizable

    \subsection{Eigenbasis}
    For a \textbf{diagonalizable} matrix $\mymat{A}$ that has distinct eigenvalues $\lambda_1, \lambda_2, \dots, \lambda_m$, \\
    \textbf{eigenbasis} is a basis $\myset{B} = \myvec{b}_1, \myvec{b}_2, \dots, \myvec{b}_n$ for the entire vector space
    \[ \operatorname{Span}(\myset{B}) = \mathbb{R}^n \]
    In other words, when the matrix is diagonalizable, the union of the bases of all its eigenspaces forms 
    an eigenbasis for the whole underlying vector space $\mathbb{R}^n$
    
    \subsection{The Diagonalization Process}
    Just as the name implies, applying a \textbf{diagonalization} on a diagonalizable matrix produces a \textbf{diagonal matrix}.
    It works as follows.
    
    \subsubsection{The way for a Matrix to record a Linear Transformation}
    Under $\mathbb{R}^n$, using the basis $\myset{B} = \{ \myvec{b}_1, \myvec{b}_2, \dots, \myvec{b}_n \}$ \vspace{1ex}

    Let $\mymat{A}$ be a $m \times n$ matrix and $L$ be a linear transformation function such that
    \[ \mymat{L}(\myvec{x}) = \mymat{A}\myvec{x} \qquad \text{ for any }n \times 1\text{ column vector }\myvec{x} \]
    
    The Matrix $\mymat{A}$ records $\mymat{L}$ by tracking how a set of vectors that equals to the basis vectors being transformed \textbf{relatively} to the basis vectors \\
    
    \begin{equation}
        \mymat{A} = \bigg[ \big[\mymat{L}(\myvec{b}_1)\big]_{\myset{B}} \, \big[\mymat{L}(\myvec{b}_2)\big]_{\myset{B}} \, \dots \, \big[\mymat{L}(\myvec{b}_n)\big]_{\myset{B}} \bigg]
    \end{equation}
    if $m \ne n$, then the $\mymat{L}$ is a transformation from $\mathbb{R}^n$ to $\mathbb{R}^m$\\
    if $m = n$, then, $\mymat{L}$ is a trasformation within the same dimension of $\mathbb{R}$
    \vspace{1ex}

    In diagonalization, we need a matrix that has eigenvectors, so let us focus on the $m = n$~case, where we have a $n \times n$~matrix~$\mymat{A}$

    \begin{minipage}{\textwidth}
        \subsubsection{Diagonalization}
        Let $\mymat{A}$ be a $n \times n$ square matrix, and $\mymat{L}$ be the linear transformation function of $\mymat{A}$ 
        \vspace{2ex}

        If we use the standard basis $\myset{S} = \{ \myvec{s}_1, \myvec{s}_2, \dots, \myvec{s}_n \}$ to express the $\mymat{L}$, in matrix, we have 
        \vspace{1ex}
    
        $[\mymat{L}(\myvec{s})]_\myset{S}$ means "how standard basis vector changes in $\mymat{L}$, relative to themselves"
        \vspace{2ex}
        
        Let's see how it becomes if we switch our basis from $\myset{S}$ to the \\
        set of eigenvectors, denote with $\myset{E} = \{ \myvec{e}_1, \myvec{e}_2, \dots, \myvec{e}_n \}$ 
        \vspace{1ex}
    
        $[\mymat{L}(\myvec{e})]_\myset{E}$ means "how eigenvector changes in $\mymat{L}$, relative to themselves", so...
      
        \[ [\mymat{L}(\myvec{e}_k)]_\myset{E} = \lambda_k [\myvec{e}_k]_{\myset{E}} = \lambda_k\,\myvec{\mymat{I}_k} \qquad \text{ where $\myvec{\mymat{I}_k}$ is the k-th column of the Identity Matrix }\]
        Therefore, 
        \begin{equation*}
            \mymat{D} = [\mymat{A}]_{\myset{E}} = \Big[ [\mymat{L}(\myvec{e}_1)]_\myset{E}\,[\mymat{L}(\myvec{e}_2)]_\myset{E}\,\dots\,[\mymat{L}(\myvec{e}_n)]_\myset{E} \Big] = \Big[ \lambda_1\myvec{\mymat{I}_1}\,\lambda_2\myvec{\mymat{I}_2}\,\dots\,\lambda_n\myvec{\mymat{I}_n} \Big]
        \end{equation*}
        $\mymat{D}$ is called the diagonalized matrix of $\mymat{A}$, by recalling Change of Basis, we have
        \begin{equation}
            \mymat{D} = \mymat{P}^{-1} \mymat{A} \mymat{P} \qquad \text{ where }\mymat{P} = \mymat{P}_{\myset{S}\leftarrow\myset{E}} = \Big[ \myvec{e}_1\,\myvec{e}_2\,\dots\,\myvec{e}_n \Big]
        \end{equation}
    \end{minipage}
    
    \subsection{Similar Matrices}
    A square matrix $\mymat{A}$ is diagonalizable if it is \textbf{similar} to a diagonal matrix $\mymat{D}$.
    This relationship is expressed through the factorization:
    \begin{equation}
        \mymat{A} = \mymat{P}\mymat{D}\mymat{P}^{-1}
    \end{equation}
    In other words, if $\mymat{D}$ is the diagonal matrix of $\mymat{A}$, $\mymat{A}$ and $\mymat{D}$ are \textbf{similar matrices} 
    \paragraph{Components of the Factorization:}
    \begin{itemize}
        \item $\mymat{D}$ is the \textbf{diagonal matrix} containing the eigenvalues $\lambda_1, \dots, \lambda_n$ along its main diagonal.
        \item $\mymat{P}$ is the \textbf{modal matrix} (or transition matrix), whose columns are the linearly independent eigenvectors $\myvec{b}_1, \dots, \myvec{b}_n$ that form the eigenbasis $\myset{B}$.
        \item $\mymat{P}^{-1}$ transforms vectors from the standard basis into the eigenbasis coordinates.
    \end{itemize}

    \paragraph{Geometric Interpretation:}
    The equation $\mymat{A} = \mymat{P}\mymat{D}\mymat{P}^{-1}$ represents a three-step process:
    \begin{enumerate}
        \item $\mymat{P}^{-1}$ changes the coordinate system of the input vector to the eigenbasis.
        \item $\mymat{D}$ performs a simple scaling (stretching/shrinking) along the axes of the eigenbasis.
        \item $\mymat{P}$ changes the resulting vector back to the standard basis.
    \end{enumerate}

    \subsection{Application: Computing Matrix Power}
    Diagonalization can optimize the computation of Matrix Power (from $\operatorname{O}(\log m)$ to $\operatorname{O}(1)$ matrix multiplications) \vspace{1ex} \\
    After we find the diagonalized matrix $\mymat{D}$ for a square matrix $\mymat{A}$, we have
    \begin{equation*}
        \mymat{D} = \mymat{P}^{-1} \mymat{A} \mymat{P} \quad\text{ or }\quad \mymat{A} = \mymat{P} \mymat{D} \mymat{P}^{-1}
    \end{equation*}
    Let us see what can we do to simplify the calculation of $\mymat{A}^m$
    \begin{align*}
        \mymat{A}^m &= (\mymat{P} \mymat{D} \mymat{P}^{-1})^m \\
        &= \overbrace{\mymat{P} \mymat{D} \mymat{P}^{-1}\,\mymat{P} \mymat{D} \mymat{P}^{-1}\,\dots\,\mymat{P} \mymat{D} \mymat{P}^{-1}}^{m\text{ times}} \\
        &= \mymat{P} \mymat{D} (\mymat{P}^{-1}\mymat{P}) \mymat{D} (\mymat{P}^{-1} \dots \mymat{P})\mymat{D}\mymat{P}^{-1} \\
        &= \mymat{P} \mymat{D}^m \mymat{P}^{-1}
    \end{align*}
    $\because$ for any diagonal matrix $\mymat{D} = \begin{bmatrix}
            d_{11}  & 0         & \dots     & 0 \\
            0       & d_{22}    & \dots     & 0 \\
            \vdots  & \vdots    & \ddots    & \vdots \\
            0       & 0         & \dots     & d_{nn}
        \end{bmatrix}$, $\mymat{D}^m = \begin{bmatrix}
            (d_{11})^m  & 0             & \dots     & 0 \\
            0           & (d_{22})^m    & \dots     & 0 \\
            \vdots      & \vdots        & \ddots    & \vdots \\
            0           & 0             & \dots     & (d_{nn})^m
        \end{bmatrix}$\\
    $\therefore$ we have optimized the calculation of $\mymat{A}^m$ from matrix multiplication of $m$ matrices \\
    to matrix multiplication of 3 matrices + scalar power of $n$ scalars

    \subsection{Orthogonal Diagonalization}
    \subsubsection{Orthonormal Set}
    \paragraph{Definition}

    An orthonormal set of vectors is basically intersecting the concept of orthogonal set and normalization of vectors
    \vspace{1ex}

    A set of vectors $\myset{W} = \{\myvec{w}_1, \myvec{w}_2, \dots, \myvec{w}_m\}$ is an \textbf{orthonormal set} \\
    if \textbf{all} the following requirements are satisfied
    \begin{itemize}
        \item $\myset{W}$ is orthogonal set
        \begin{itemize}
            \item $\dotprod{\myvec{w}_i}{\myvec{w}_j} = 0 \quad \forall\,i, j \in \{1, \dots, m\} \text{ with } i \neq j$
        \end{itemize}
        \item $\|\myvec{w}_k\| = 1 \quad \forall\,k \in \{1, \dots, m\}$
    \end{itemize}
    \paragraph{Note:} $\myset{W} \not\subset \mathbb{R}^n$

    
    \subsubsection{Orthogonal Complements}
    \paragraph{Orthogonal to a Set of Vectors}
    \begin{quote}
        a vector $\myvec{z}$ is orthogonal to a vector space $\mysubvecsp{V}$ if \\
        \[ \dotprod{\myvec{z}}{\myvec{v}} = 0 \quad \forall\,\myvec{v} \in \mysubvecsp{V} \]
    \end{quote}
    
    \paragraph{Orthogonal Complement of a Set of Vectors} 
    
    \begin{quote}
        The \textbf{orthogonal complement} of a vector space $\mysubvecsp{V}$ (denoted~by~$\mysubvecsp{V}^{\perp}$) is the set of \textbf{all}~vectors~$\myvec{z}$ that are orthogonal to $\mysubvecsp{V}$
        \vspace{1ex}
        Note:
        \begin{itemize}
            \item $\mysubvecsp{V}^{\perp} \subset \mathbb{R}^n$
            \item any $\myvec{z} \in \mysubvecsp{V}^{\perp}$ is also orthogonal to $\operatorname{Span}(\mysubvecsp{V})$, $\implies \mysubvecsp{V}^{\perp} = (\operatorname{Span}(\mysubvecsp{V}))^{\perp}$
        \end{itemize}
    \end{quote}
    
    \newpage
    \subsubsection{Matrices with Orthonormal Column} \label{section3_5_3}
    A $m \times n$ matrix with orthonormal columns implies that \\
    by taking the columns as column vectors, all the column vectors forms an orthogonal set
    
    \paragraph{$\bm{\mymat{U}^T \mymat{U} = \mymat{I}}$ :}\mbox{}\vspace{1ex}

    Let $\mymat{U} = \begin{bmatrix} \myvec{u}_1 & \myvec{u}_2 & \dots & \myvec{u}_n \end{bmatrix}$ be a $m \times n$ matrix, then
    \begin{equation}
        \mymat{U}\text{ has orthonormal columns} \iff \mymat{U}^T \mymat{U} = \mymat{I}
    \end{equation}
    Proof:
    \begin{quote}
        \begin{align*}
            \renewcommand{\arraystretch}{1.5}
            \mymat{U}^{T} \mymat{U} = 
                \begin{bmatrix} \myvec{u}_1^T \\ \myvec{u}_2^T \\ \vdots \\ \myvec{u}_n^T \end{bmatrix}
                \begin{bmatrix} \myvec{u}_1 & \myvec{u}_2 & \dots & \myvec{u}_n \end{bmatrix}
                &= \begin{bmatrix}
                    \dotprod{\myvec{u}_1}{\myvec{u}_1} & \dotprod{\myvec{u}_1}{\myvec{u}_2} & \dots  & \dotprod{\myvec{u}_1}{\myvec{u}_n} \\
                    \dotprod{\myvec{u}_2}{\myvec{u}_1} & \dotprod{\myvec{u}_2}{\myvec{u}_2} & \dots  & \dotprod{\myvec{u}_2}{\myvec{u}_n} \\
                    \vdots & \vdots & \ddots & \vdots \\
                    \dotprod{\myvec{u}_n}{\myvec{u}_1} & \dotprod{\myvec{u}_n}{\myvec{u}_2} & \dots  & \dotprod{\myvec{u}_n}{\myvec{u}_n}
                \end{bmatrix}
        \end{align*}
        Note that if $\mymat{U}$ has orthonormal column, by definition of orthonormal set
        \begin{equation*}
            \dotprod{\myvec{u}_i}{\myvec{u}_j} = \delta_{ij} = \begin{cases}
                1 & \text{if } i = j \\
                0 & \text{if } i \ne j
            \end{cases} \quad \forall\,i, j \in \{1, \dots, n\}
        \end{equation*}
        Therefore,
        \begin{equation*}
            \mymat{U}^T\mymat{U} = \begin{bmatrix}
                1       & 0         & \dots     & 0 \\
                0       & 1         & \dots     & 0 \\
                \vdots  & \vdots    & \ddots    & \vdots \\
                0       & 0         & \dots     & 1 
            \end{bmatrix} = \mymat{I}_{n \times n}
        \end{equation*}
    \end{quote}

    \paragraph{Other Properties}\mbox{}\vspace{1ex}

    Let $\mymat{U} = \begin{bmatrix} \myvec{u}_1 & \myvec{u}_2 & \dots & \myvec{u}_n \end{bmatrix}$ be a $m \times n$ matrix with orthonormal columns \\
    and let $\myvec{x}, \myvec{y} \in \mathbb{R}^n$
    \begin{itemize}
        \item \textbf{Length Preservation:} $\|\mymat{U}\myvec{x}\| = \|\myvec{x}\|$ for all $\myvec{x} \in \mathbb{R}^n$.
        \item \textbf{Dot Product Preservation:} $(\mymat{U}\myvec{x}) \cdot (\mymat{U}\myvec{y}) = \myvec{x} \cdot \myvec{y}$ for all $\myvec{x}, \myvec{y} \in \mathbb{R}^n$.
        \item $\dotprod{\myvec{x}}{\myvec{y}} = 0 \iff \dotprod{(\mymat{U}\myvec{x})}{(\mymat{U}\myvec{y})} = 0$            
    \end{itemize}
    
    \newpage
    \subsubsection{Orthogonal Matrix}
    \paragraph{Definition: } An orthogonal matrix is an invertible (square) matrix $\mymat{U}$ such that
    \begin{equation}
        \mymat{U}^{-1} = \mymat{U}^T
    \end{equation}
    \paragraph{Notable Properties: }
    \begin{itemize}
        \item $\mymat{U}$ is an orthogonal matrix \\ 
            $\iff \mymat{U}$ is a square matrix with ortho\textbf{normal} columns \\
            \quad $\implies \mymat{U}^T \mymat{U} = \mymat{I}$ (see subsubsection \ref{section3_5_3})\\
            $\iff \mymat{U}$ is a square matrix with ortho\textbf{normal} rows \\
            \quad $\implies \mymat{U} \mymat{U}^T = \mymat{I}$ \\
        \item The determinant of an orthogonal matrix is always $\pm 1$ \\
            Proof: 
            \begin{quote}
                Let $\mymat{U}$ be any orthogonal matrix, then $\det \mymat{U} \ne 0$,
                \begin{align*}
                    \mymat{U}^{-1} &= \mymat{U}^T \\
                    \det(\mymat{U}^{-1}) &= \det(\mymat{U}^T) \\
                    \frac{1}{\det\mymat{U}} &= \det \mymat{U} \\
                    ( \det \mymat{U} )^2 &= 1 \\
                    \Aboxed{\det \mymat{U} &= \pm 1}
                \end{align*}
            \end{quote}
        \end{itemize}
    Because of some historic reason, such a matrix is called an orthogonal matrix but not an orthonormal matrix, 
    but an orthogonal matrix actually \textbf{requires} its set of column vectors (or set of row vectors) to be \textbf{orthogonal and normalized}.
    
    \subsubsection{Symmetric Matrix}
    A matrix $\mymat{A}$ is a symmetric matrix if 
    \begin{equation}
        \mymat{A}^T = \mymat{A}
    \end{equation}
    A symmetric matrix has its entries symmetric along the diagonal from top-left to bottom-right

    \newpage
    \subsubsection{Orthogonally Diagonalizable} \label{sec3_5_6}
    \paragraph{Definition} An $n \times n$ matrix $\mymat{A}$ is \textbf{orthogonally diagonalizable} if there are an orthogonal matrix $\mymat{P}$ and a diagonal matrix $\mymat{D}$ such that
    \begin{equation}
        \mymat{A} = \mymat{P}\mymat{D}\mymat{P}^{T} = \mymat{P}\mymat{D}\mymat{P}^{-1}
    \end{equation}
    Similar to Orthogonal Matrix, the naming is a bit confusing. \\
    A \textbf{orthogonally diagonalizable} matrix means that \\
    \quad the matrix $\mymat{A}$ can be \textbf{diagonalized} by changing the basis into a \\
    \quad set of \textbf{orthogonal and normalized}~basis~vector~(namely~ortho\textbf{normal}~basis)
    Therefore, the transition matrix $\mymat{P}$ has to be \\
    \quad an orthogonal matrix (with ortho\textbf{normal} columns) in \\
    \quad order for matrix to be orthogonally diagonalizable

    \paragraph{The Spectral Theorem for Real Symmetric Matrices}~\\
    \quad Let $\mymat{A}$ be a $n \times n$ orthogonally diagonalizable matrix
    \begin{enumerate}
        \item orthogonally diagonalizable matrix $\iff$ \textbf{symmetric} matrix ($\mymat{A}^T = \mymat{A}$)
        \item all eigenvalues of $\mymat{A}$ are \textbf{real}
        \item for each eigenvalue $\lambda_k$, geometric multiplicity = algebraic multiplicity
        \item eigenspaces of different eigenvalues are \textbf{orthogonal complements} of each other
    \end{enumerate}

    \paragraph{Proofs of Properties}
    \noindent
    \begin{minipage}{\textwidth}
        1. orthogonally diagonalizable matrix $\iff$ \textbf{symmetric} matrix ($\mymat{A}^T = \mymat{A}$)
        \begin{quote}
            
        \end{quote}
    \end{minipage} 
    \vspace{2ex}

    \begin{minipage}{\textwidth}
        2. all eigenvalues of a real symmetrix matrix $\mymat{A}$ are \textbf{real}
        \begin{quote}
            Let $\lambda_k$ be any eigenvalue of the real matrix $\mymat{A}$ , \\ 
            and $\myvec{v}_k$ be the corresponding eigenvector (potentially complex)
            \begin{align*}
                \mymat{A}\myvec{v}_k &= \lambda_k\myvec{v}_k \\
                (\mymat{A}\myvec{v}_k)^{*} &= (\lambda_k\myvec{v}_k)^{*} & \text{take conjugate transpose} \\
                (\myvec{v}_k)^{*} \mymat{A}^T &= \bar{\lambda_k}\,(\myvec{v}_k)^{*} & \text{properties of conjugate transpose} \\
                (\myvec{v}_k)^{*} \mymat{A} &= \bar{\lambda_k}\,(\myvec{v}_k)^{*} & \text{$\mymat{A}$ is similar} \\
                (\myvec{v}_k)^{*} \mymat{A}^T\,\myvec{v}_k &= \bar{\lambda_k}\,(\myvec{v}_k)^{*}\,\myvec{v}_k & \text{right-multiply by }\myvec{v}_k \\
                (\myvec{v}_k)^{*} (\lambda_k \myvec{v}_k) &= \bar{\lambda_k}\,(\myvec{v}_k)^{*}\,\myvec{v}_k & \\
                \lambda_k (\myvec{v}_k)^{*} \myvec{v}_k &= \bar{\lambda_k}\,(\myvec{v}_k)^{*}\,\myvec{v}_k & \\
                \lambda_k &= \bar{\lambda_k} & (\myvec{v}_k)^{*}\myvec{v}_k = \|\myvec{v}\|^2 \ne 0 \\
                \Aboxed{\lambda_k &\in \mathbb{R}}
            \end{align*}
        \end{quote}    
    \end{minipage} 
    \vspace{2ex}

    \noindent
    \begin{minipage}{\textwidth}
        4. eigenspaces of different eigenvalues are \textbf{orthogonal complements} of each other
        \begin{quote}
            Let $\lambda_1$ and $\lambda_2$ be two distinct eigenvalues of $\mymat{A}$ \\
            and let $\myvec{v}_1$ and $\myvec{v}_2$ be the eigenvectors corresponds to $\lambda_1$ and $\lambda_2$ respectively
            \begin{align*}
                \lambda_1 \dotprod{\myvec{v}_1}{\myvec{v}_2} &= \dotprod{(\lambda_1 \myvec{v}_1)}{\myvec{v}_2} \\
                &= \dotprod{(\mymat{A} \myvec{v}_1)}{\myvec{v}_2} \\
                &= (\mymat{A} \myvec{v}_1)^T \myvec{v}_2 & \dotprod{\myvec{u}_1}{\myvec{u}_2} = \myvec{u}_1^T \myvec{u}_2 \text{ for all real column vectors } \\
                &= (\myvec{v}_1)^T \mymat{A}^T \myvec{v}_2 \\
                &= (\myvec{v}_1)^T \mymat{A} \myvec{v}_2 & \mymat{A}^T = \mymat{A}\\
                &= (\myvec{v}_1)^T (\lambda_2 \myvec{v}_2) \\
                &= \lambda_2 \dotprod{\myvec{v}_1}{\myvec{v}_2} \\
                \lambda_1 \dotprod{\myvec{v}_1}{\myvec{v}_2} &= \lambda_2 \dotprod{\myvec{v}_1}{\myvec{v}_2} \\
                (\lambda_1 - \lambda_2) \dotprod{\myvec{v}_1}{\myvec{v}_2} &= 0 
            \end{align*}
            Since $\lambda_1 \ne \lambda_2$, we have $\dotprod{\myvec{v}_1}{\myvec{v}_2} = 0$ \\ 
            Therefore, for any two distinct eigenvalues of a symmetric matrix, the eigenvectors corresponds to each of them are perpendicular to each other
        \end{quote}
    \end{minipage}

    \newpage
    \subsubsection{Orthogonal Diagonalization of Symmetric Matrices}
    \paragraph{Gram-Schmidt Orthogonalization Process}~\vspace{1ex}

    The Gram-Schmidt Orthogonalization Process is a useful algorithmic procedure.
    \\
    Take it briefly, given a set of basis vectors of any vector space $\mysubvecsp{V} \subseteq \mathbb{R}^n$,\\
    \quad the process find a set of orthogonal (but not normalized) basis for $\mysubvecsp{V}$ \vspace{1ex}
    
    It is as follows:
    \begin{quote}
        Let $\{\myvec{w}_1, \myvec{w}_2, \dots, \myvec{w}_n\}$ be a basis for any vector space $\mysubvecsp{V} \subseteq \mathbb{R}^n$, then\\
        the orthogonal basis $\{\myvec{v}_1, \myvec{v}_2, \dots, \myvec{v}_n\}$ for $\mysubvecsp{V}$ can be found as follows:
    
        \begin{align*}
            \myvec{v}_1 &= \myvec{w}_1 \\
            \myvec{v}_2 &= \myvec{w}_2 - \frac{\dotprod{\myvec{v}_1}{\myvec{w}_2}}{\dotprod{\myvec{v}_1}{\myvec{v}_1}}\myvec{v}_1 \\
            \myvec{v}_3 &= \myvec{w}_3 - \frac{\dotprod{\myvec{v}_1}{\myvec{w}_3}}{\dotprod{\myvec{v}_1}{\myvec{v}_1}}\myvec{v}_1 - \frac{\dotprod{\myvec{v}_2}{\myvec{w}_3}}{\dotprod{\myvec{v}_2}{\myvec{v}_2}}\myvec{v}_2 \\
            \dots \\
            \myvec{v}_n &= \myvec{w}_n - \sum_{k = 1}^{n-1} \frac{\dotprod{\myvec{v}_k}{\myvec{w}_n}}{\dotprod{\myvec{v}_k}{\myvec{v}_k}}\myvec{v}_k \\
        \end{align*}
    \end{quote}
    The Gram-Schmidt Orthogonalization Process is useful in finding an orthogonal diagonalization.
    
    \begin{minipage}{\textwidth}
        \raggedright
        \paragraph{Process of Orthogonal Diagonalization}~\vspace{1ex}\\
        Given a $n \times n$ square matrix $\mymat{A}$, if we want to diagonalize it orthogonally, the steps are:
        \begin{enumerate}
            \item check if $\mymat{A}$ is orthogonally diagonalizable, i.e. a symmetric matrix
            \item find all distinct eigenvalues $\lambda_1, \lambda_2, \dots, \lambda_m$ from $\det(\mymat{A} - \lambda\mymat{I}) = 0$
            \item substitute all $\lambda_k$ one by one and solve the equation $\mymat{A}\myvec{x}_k=\lambda_k\myvec{x}_k$ for $\myvec{x}_k$
            \item[-.] we should now have $\myvec{x}_1, \myvec{x}_2, \dots, \myvec{x}_m$, \\
                some of them may have 1 free variable only, while others have more free variables    
            \item if $\myvec{x}_k$ has only 1 free variable, because of property 4 mentioned in \ref{sec3_5_6}, it can directly be a vector in the orthogonal basis $\myvec{u}$\vspace{1ex} \\
                if $\myvec{x}_k$ has more (num = $l$) free variables, get the basis vectors $\myvec{w}_1, \myvec{w}_2, \dots, \myvec{w}_l$ for the eigenspace of $\myvec{x}_k$, 
                and then use the Gram-Schmidt Orthogonalization Process to process the $\myvec{w}$s into orthogonal basis vectors $\myvec{u}_1, \myvec{u}_2, \dots, \myvec{u}_l$
            \item[-.] Collect the resulting orthogonal vector $\myvec{u}$s from all the $\myvec{x}$s, there should be $n$ of them, they form the orthogonal eigenbasis
            \item normalize all $\myvec{u}$s of the orthogonal eigenbasis into orthnormal eigenbasis $\{ \myvec{v}_1, \myvec{v}_2, \dots, \myvec{v}_n \}$
            \item put the $\myvec{v}$s as columns into a matrix, this is the $\mymat{P}_{\myset{S}\leftarrow\myset{E}}$
            \item put the eigenvalues $\lambda_1, \dots, \lambda_m$ in the same order as the vector $\myvec{v}$s (a $\lambda_k$ may corresponds to multiple $\myvec{v}$s, just put duplicated $\lambda$s as if u treat the duplicated $\lambda$s differently from the characteristic equation) into the diagonals into a $n \times n$ matrix, this is the diagonalized matrix $\mymat{D}$
        \end{enumerate}
    \end{minipage}
    
    \newpage
    \vspace*{\fill}
    \begin{center}
        \LARGE
        \fontfamily{qzc}\selectfont\itshape the end
    \end{center}
    \vspace*{\fill}
    
\end{document}